\begin{table*}[!tb]
  \centering
  \caption{The 100 evaluated proteins against the full per-organism dataset that they were randomly drawn from. Protein length and fragment length are in residues.}
  \label{tab:dataset_coverage}
  \begin{tabular}{lrrrr}
    \toprule
    Group & P. Count & P. Length Mean & F. Count Mean & F. Length Mean \\
    \midrule
    \textit{E. coli} - Full Dataset & 5118 & 309.4 & 19.6 & 15.8 \\
    \textit{E. coli} - Evaluated Split & 100 & 319.5 & 19.3 & 16.6 \\
    \midrule
    \textit{S. cerevisiae} - Full Dataset & 6713 & 452.6 & 31.8 & 14.2 \\
    \textit{S. cerevisiae} - Evaluated Split & 100 & 442.3 & 32.5 & 13.6 \\
    \bottomrule
  \end{tabular}
  \par\smallskip\footnotesize{P.\ = protein, F.\ = fragment; lengths are in amino acids (aa). F.\ Count Mean is the mean number of fragments a protein digests into. Both groups are measured on the same digestion replica, so any difference is sampling rather than digestion randomness.}
\end{table*}
